\documentclass[10pt,a4paper]{ujarticle}
\usepackage[utf8]{inputenc}
\usepackage{amsmath}
\usepackage{amsfonts}
\usepackage{amssymb}
\usepackage{problem}
\usepackage{amsthm}
\newtheorem{theorem}{定理}
\newtheorem{lemma}{補題}
\newtheorem{definition}{定義}
\usepackage {comment}
\usepackage{bm}
\author{山田 航輝}
\usepackage[top=30truemm,bottom=30truemm,left=25truemm,right=25truemm]{geometry}
\usepackage{multicol}
\renewcommand{\labelenumi}{(\arabic{enumi})}
\title{LagrangeスペクトラムとHall's ray}
\begin{document}
\maketitle
\begin{abstract}
本稿では，Lagrange スペクトラムの定義，およびLagrange スペクトラムがHall's rayと呼ばれる半直線を含むことの証明を述べる．また，それらがどのように形式化されているかについても記述する．
\end{abstract}

\section{Lagrange スペクトラムとHall's ray}
まず．本レポジトリで形式化するLagrangeスペクトラムの定義と，Hall's rayの存在について述べる．
\begin{definition}
自然数の両側無限列$(a_n)_{n\in\mathbb{Z}}\in (\mathbb{N}_{\ge 1})^\mathbb{Z}$と整数$i$に対して，
$$
\lambda_i((a_n)_{n\in\mathbb{Z}}) = [a_i; a_{i+1}, a_{i+2},\cdots] + [0; a_{i-1}, a_{i-2},\cdots]
$$
と定義する．ただし，$[c_0; c_1, c_2, \cdots]$は連分数
$$
c_0 + \frac{1}{c_1 + \frac{1}{c_2 + \cdots}}
$$
を表す．


\textbf{Lagrangeスペクトラム}とは，
$$
L = \left\{\limsup_{i \to \infty}\lambda_i ((a_n)_{n\in\mathbb{Z}}) : (a_n)_{n\in\mathbb{Z}}\in (\mathbb{N}_{\ge 1})^\mathbb{Z}, \limsup_{i \to \infty}\lambda_i ((a_n)_{n\in\mathbb{Z}}) < \infty \right\}
$$
で定義される実数の部分集合である．
\end{definition}
\begin{theorem}[Hall, 1947]
\label{theorem: C4}
$$
C(N) = \{[0;a_1, a_2, \cdots] : 1\le a_i \le N\}
$$
とする．このとき，
$$
C(4) + C(4) = \{x + y : x\in C(4), y\in C(4)\} = [\sqrt{2} - 1, 4\sqrt{2} - 4]
$$
である．
\end{theorem}
この系として，

\begin{theorem}
\label{theorem: ray}
$$
[6,\infty) \subset L.
$$
\end{theorem}
が得られる．これと$L$が閉集合であること（これはまだ形式化していない）から，ある実数$c_F$より大きい実数が$L$に含まれることが分かる．$L$に含まれるこの半直線$[c_F, \infty)$を\textbf{Hall's ray}と呼ぶ．

\section{Hall's rayが存在することの証明}
\subsection{補題およびその証明}
\cite{cusick}に従って証明する．
\begin{lemma}
\label{lemma: remove_interval}
$a_1,\cdots, a_n \in \mathbb{N}_{\ge 1}$とし，以下の三つの型の区間を考える．
\begin{enumerate}
\item $[0;a_1,\cdots, a_n ,\overline{1,4}]$と$[0;a_1,\cdots, a_n ,\overline{4,1}]$を端点とする区間
\item $[0;a_1,\cdots, a_n ,\overline{4,1}]$と$[0;a_1,\cdots, a_n , 2, \overline{4,1}]$を端点とする区間
\item $[0;a_1,\cdots, a_n ,\overline{4,1}]$と$[0;a_1,\cdots, a_n , 3, \overline{4,1}]$を端点とする区間
\end{enumerate}
そして，
\begin{enumerate}
\item からは$[0;a_1,\cdots, a_n , 2, \overline{4,1}]$と$[0;a_1,\cdots, a_n , 1, \overline{1,4}]$を端点とする区間を
\item からは$[0;a_1,\cdots, a_n , 3, \overline{4,1}]$と$[0;a_1,\cdots, a_n , 2, \overline{1,4}]$を端点とする区間を
\item からは$[0;a_1,\cdots, a_n , 4, \overline{4,1}]$と$[0;a_1,\cdots, a_n , 3, \overline{1,4}]$を端点とする区間を
\end{enumerate}
それぞれ取り除くことを考える．このプロセスによって
\begin{enumerate}
\item の型の区間からは$r_{n+1} = 1$の（1）と（2）の型の区間が，
\item の型の区間からは$r_{n+1} = 1$の（1）と（3）の型の区間が，
\item の型の区間からは$r_{n+1} = 3$の（1）と$r_{n+1} = 4$の（1）の型の区間が
\end{enumerate}
それぞれ得られる．そして，このプロセスによって除かれる区間を$I$とし，生じる区間を$M_1, M_2$とすると，
$$
|M_i|\ge |I|
$$である．ただし$|A|$は区間$A$の長さを表す．
\end{lemma}
\begin{proof}
あとで書く
\end{proof}

\begin{lemma}
\label{lemma: two_closed}
$A_1, A_2$を有界な閉区間とし，$A_1$から開区間$I$を取り除いて左から順に$A_{11}, A_{12}$が生じたとする．このとき，$|A_2| \ge |I|$なら，
$$
A_1 + A_2 = (A_{11} \cup A_{12}) + A_2.
$$
\end{lemma}
\begin{proof}
$A_{11} + A_2$は$A_1 + A_2$と左端の点$l$を共有する．また，$A_{11} + A_2$の右端の点は$l + |A_{11}| + |A_2|$である．

一方，$A_{12} + A_2$は$A_1 + A_2$と右端を共有し，左端の点は$l + |A_{11}| + |I|$である．$|A_2| \ge |I|$より，
$A_{12} + A_2$の左端の点の方が$A_{11} + A_2$の右端の点よりも小さくなる．よって$(A_{11} \cup A_{12}) + A_2$は連結な閉区間であり，$A_1 + A_2$の右端と左端の点を含むから，$A_1 + A_2 = (A_{11} \cup A_{12}) + A_2$である．
\end{proof}
\begin{lemma}
\label{lemma: sum_invariance}
$A_1, A_2,\cdots A_r$を互いに素な有界閉区間の列とし，$A_1$から開区間$I$を取り除いて左から順に$A_{r+1}, A_{r+2}$が生じたとする．
$$
B = \bigcup_{i=1}^r A_i,\quad B^* = \bigcup_{i=2}^{r+2} A_i
$$
とする．このとき，すべての$2\le i\le r+2$について$|A_i| \ge |I|$なら，$B^* + B^* = B + B$
\end{lemma}
\begin{proof}
$$
A_1^* = A_{r+1} \cup A_{r+2}, \quad A_i^* = A_i\  (2\le i \le r)
$$とする．
補題\ref{lemma: two_closed}から，
\begin{align*}
A_1 + A_1 &= (A_{r+1} \cup A_{r+2}) + A_1 \\
&=(A_{r+1} + A_1) \cup (A_{r+2} + A_1) \\
&=(A_{r+1} + (A_{r+1} \cup A_{r+2})) \cup (A_{r+2} + (A_{r+1} \cup A_{r+2})) \\
&=A_1^* + A_1^*, \\
A_1 + A_i &= A_1^* + A_i^* \ (2\le i \le r)
\end{align*}
が分かる．よって
\begin{align*}
B + B &= \left(\bigcup_{j = 1}^r (A_1 + A_j)\right) + \left(\bigcup_{2\le i \le j \le r} (A_i + A_j)\right)\\
&= \left(\bigcup_{j = 1}^r (A_1^* + A_j^*)\right) + \left(\bigcup_{2\le i \le j \le r} (A_i^* + A_j^*)\right)\\
&= \bigcup_{1\le i \le j \le r} (A_i^* + A_j^*)\\
&= B^* + B^*
\end{align*}
\end{proof}
\begin{lemma}
\label{lemma: sum_cap_eq_cap_sum}
$C_1,C_2,\cdots$を有界閉集合の列とし，$C_{i+1} \subset C_i$とする．このとき，
$$
\bigcap_{i = 1}^\infty C_i + \bigcap_{i = 1}^\infty C_i = \bigcap_{i = 1}^\infty (C_i + C_i).
$$
\end{lemma}
\begin{proof}
各$i$について，$C_i + C_i$は左辺を含むから，$（左辺）\subset （右辺）$はすぐに分かる．

$t\in \bigcap_{i = 1}^\infty (C_i + C_i)$とする．そして各$i$について，$V_i = \{(x_1,x_2) : x_1, x_2 \in C_i \}$とし，$T = \{(x_1,x_2) : x_1+ x_2 = t \}$とする．$V_i$は閉集合で有界，$T$は閉集合であるから，$T\cap V_i$は有界閉集合である．$t\in C_i + C_i$だから，$T\cap V_i \neq \emptyset$である．$T\cap V_{i+1} \subset T \cap V_i$なので，
$$
\bigcap_{i=1}^\infty (T \cap V_i) = T\cap \bigcap_{i=1}^\infty V_i \neq \emptyset
$$
である．よって，ある$(x_1, x_2) \in T\cap \bigcap_{i=1}^\infty V_i \neq \emptyset$がとれる．$(x_1, x_2)\in T$だから$x_1 + x_2 = t$，また$(x_1, x_2)\in \bigcap_{i=1}^\infty V_i$より$x_1, x_2 \in \bigcap_{i=1}^\infty C_i$となり，
$$
\bigcap_{i = 1}^\infty (C_i + C_i) \subset 
\bigcap_{i = 1}^\infty C_i + \bigcap_{i = 1}^\infty C_i
$$
が分かる．以上より，
$$
\bigcap_{i = 1}^\infty C_i + \bigcap_{i = 1}^\infty C_i = \bigcap_{i = 1}^\infty (C_i + C_i).
$$
となる．
\end{proof}

\subsection{定理\ref{theorem: C4}の証明}
まず，
$S_0 = [(\sqrt{2} - 1)/2, 2(\sqrt{2} - 1)]$とし，
\begin{itemize}
\item $S_i \supset S_{i + 1}$
\item $\bigcap_{i=0}^\infty S_i = C(4)$
\item $S_i + S_i = S_{i + 1} + S_{i + 1}$
\end{itemize}
となる有界閉集合の列を構成する．

$$
[0;\overline{4, 1}] = \frac{\sqrt{2} - 1}{2}, \quad [0;\overline{1, 4}] = 2(\sqrt{2} - 1)
$$
であり，これらは$C(4)$の最小と最大の元である．また，$S_0$は補題\ref{lemma: remove_interval}における（1）の型の区間である．

$C(4)$は区間$S_0$から，補題\ref{lemma: remove_interval}の手順に従って（可算無限個）区間を取り除いていくことで得られる．ここでは，順番を変えて，長い順に取り除くことにし，$D_0, D_1, D_2, \cdots$の順で取り除くことにする．したがって，$|D_i| \ge |D_{i +1}|$である．

そして，$S_i = S_{i-1}\setminus D_{i -1}$とする．$B = S_i$，$B^* = S_{i + 1}$，$I = D_i$に対して補題\ref{lemma: sum_invariance}を適用するために，$D_i$が$S_{i + 1}$に残るどの区間よりも小さいことを示す．

$S_i$内の閉区間$H_i$から$D_i$を取り除くことで，$S_{i+1}$が得られるとしよう．$D_i$をとり除くことで，$H_i$は$H_{i1}, H_{i2}$の二つの区間に分かれる．ここで，$B = S_{i - 1}$，$B^* = S_{i}$，$I = D_{i- 1}$に対して補題\ref{lemma: sum_invariance}の適用条件が満たされていたとすれば，$|D_{i - 1}| \ge |D_i|$であるから，閉区間に関する$|H_{ij}| \ge |D_i|$を示せば十分である．

補題\ref{lemma: remove_interval}の手順に従って区間を取り除いていったとき，$D_i$が取り除かれる際にできる閉区間を$M_1, M_2$とする．補題\ref{lemma: remove_interval}より，$|M_j| \ge |D_i|$であり，$H_{ij}$が$M_1$または$M_2$と等しい，または$M_1, M_2$を含めば，$|H_{ij}| \ge |D_i|$である，

$H_{ij}$が$M_1$または$M_2$に真に含まれていたとすると，$H_{ij}$の$D_i$とは反対側で取り除かれる区間$D_m$は$M_1$ないし$M_2$に含まれる．したがって，補題\ref{lemma: remove_interval}の手順に従って区間を取り除いていったとき，$D_m$は，$D_i$の後に取り除かれなくてはならない．よって，その際$D_m$を取り除くことによって得られる区間$N$（$D_m$と$D_i$に挟まれている側）は補題\ref{lemma: remove_interval}から$|N|\ge |D_m|$を満たす．また，大きい順に取り除いていったとき，$D_m$は$D_i$より先に取り除かれていたのだから，$m < i$であり，
$$
|H_{ij}| \ge |N| \ge |D_m| \ge |D_i|
$$
が分かる．よって$B = S_i$，$B^* = S_{i + 1}$，$I = D_i$に対して補題\ref{lemma: sum_invariance}を適用でき，$S_i + S_i = S_{i + 1} + S_{i + 1}$となる．作り方から$S_i supset S_{i + 1}$や$\bigcap_{i=0}^\infty S_i = C(4)$だから，求める有界閉集合の列が得られた．

よって補題\ref{lemma: sum_cap_eq_cap_sum}から，
\begin{align*}
C(4) + C(4) &= \bigcap_{i = 0}^\infty S_i + \bigcap_{i = 0}^\infty S_i \\
&= \bigcap_{i = 0}^\infty (S_i + S_i) \\
&= S_0 + S_0\\
&= [\sqrt{2} - 1, 4\sqrt{2} - 4]
\end{align*}
分かる．\hfill \qed
\subsection{定理\ref{theorem: C4}による定理\ref{theorem: ray}の証明}
$[\sqrt{2} - 1, 4\sqrt{2} - 4]$の大きさは$1$より大きいから，$6\le l$をとると，
$$
l = a + [0;b_1, b_2, \cdots ] + [0;c_1, c_2, \cdots ]
$$
と表せる．ここで，$a$は$5$以上の整数，$b_i, c_i$は$1\le b_i,c_i\le 4$なる整数である．

自然数の両側無限列$A$を
$$
A = (\cdots, 1, 1, a, b_1, a, c_1, b_2, b_1, a, c_1, c_2,  b_3, b_2, b_1, a, c_1, c_2, c_3, \cdots)
$$
と定める（最初の$a$の場所を$0$番目としておく）．すると，$a$が$5$以上だから，それ以外の場所での$\lambda_i$の値は$\limsup$を取る際に無視してよく，
$$
\limsup_{i \to \infty} \lambda_i(A) = l
$$
となる．つまり$l \in L$となる．したがって，
$$
[6,\infty) \subset L
$$
が分かる．\hfill \qed
\begin{thebibliography}{99}
%\bibitem{ma79} Markoff, A. (1879) {\it Sur les formes quadratiques binaires indéfinies.} Math. Ann. 15, 381–406. 
%\bibitem{ma80} Markoff, A. (1880) {\it Sur les formes quadratiques binaires indéfinies.} Math. Ann. 17, 379–399. 
%\bibitem{hall} Hall, M. (1947). On the Sum and Product of Continued Fractions. Annals of Mathematics, 48(4), 966–993. 
\bibitem{cusick} T.~W. Cusick and M.~E. Flahive, {\it The Markoff and Lagrange spectra}, Mathematical Surveys and Monographs, 30, Amer. Math. Soc., Providence, RI, 1989; MR1010419
\end{thebibliography}
\end{document}